\section{Discussion and conclusion}
\label{sec:discussion}
A distribution repair needs a target that survives changes in how
samples are pooled and compared. Fixing the reference and a meaningful
effect size makes regional probability loss testable. Conditional spread
and fidelity then show whether changes in variation serve that target.
The crossed probes illustrate the distinction: longer sampling improves
some short-sampler deficits, while the missing digit persists across
adequate samplers when training withholds it.

The two controlled tests establish partial repairs at matched budgets.
Starting-covariance matching improves two-call agreement with real data
on one pretrained model, even though most teacher-defined regions remain
depleted. Ordinary continuation improves the full distribution of digit
classifier outcomes across three paired seeds, including rejection in
the measure. Its rejection rate still exceeds that of real data. A second
pretrained checkpoint would test transfer of the initialization result,
and independent labels for generated digits would test whether the
classifier findings track human judgments. The matched-cost prompt
shuffle leaves the advantage of per-prompt selection unresolved.
